% ===========================================================================
%  控制设计约束与 Web 基线
% ===========================================================================

\section{控制设计约束与 Web 基线控制器}
\label{sec:control}

\subsection{两套控制实现的角色}

工程中同时存在 Web 交互仿真控制器和飞控固件控制器。Web 控制器用于构型可视化、
确定性回归和概念交互；固件控制器用于 200\,Hz 实时闭环。二者共享推进方程和部分参数，
但当前并非同一算法：Web 垂直模式使用
$\bm q_d^{-1}\otimes\bm q$ 的误差方向和累加式局部分配，固件使用
$\bm q^{-1}\otimes\bm q_d$ 和绝对力矩分配。本文因此只把本节算法称为“Web 基线”，
不以其结果替代\cref{sec:cascade_chapter}的固件验证。

\subsection{前飞基线 SAS}

Web 前飞模式在推进模型系中使用角速率阻尼与两个姿态保持项。令
$p',q',r'$ 表示该模型系的角速度，则
\begin{align}
\delta_t^{\mathrm{cmd}}
&=\delta_{t0}+k_q q'-k_\theta(\theta-\theta_0)-k_{I\theta}\!\int(\theta-\theta_0)\,\mathrm dt,
\label{eq:sas_pitch}\\
\delta_f^{\mathrm{cmd}}
&=\delta_{f0}-k_r r',
\label{eq:sas_yaw}\\
\Delta\omega^{\mathrm{cmd}}
&=k_p p'-k_\phi\phi-k_{I\phi}\!\int\phi\,\mathrm dt.
\label{eq:sas_roll}
\end{align}
式\eqref{eq:sas_pitch}--\eqref{eq:sas_roll}的符号由
$\partial M_{y'}/\partial\delta_t<0$、
$\partial M_{z'}/\partial\delta_f>0$ 和
$\partial M_{x'}/\partial\Delta\omega<0$ 决定。
这里的变量名称沿用 Web 模型轴，不能直接替换成固件悬停滚转/俯仰/偏航名称。

角速度命令模式把用户输入解释为模型角速度参考：
\begin{equation}
\begin{aligned}
\delta_t^{\mathrm{cmd}}&=\delta_{t0}+k_q(q'-q'_{\mathrm{ref}}),\\
\delta_f^{\mathrm{cmd}}&=\delta_{f0}+k_r(r'_{\mathrm{ref}}-r'),\\
\Delta\omega^{\mathrm{cmd}}&=k_p(p'-p'_{\mathrm{ref}}).
\end{aligned}
\label{eq:rate_command}
\end{equation}
该控制器主要用于检查通道方向、限幅和回归样本，不包含固件的 EKF 零偏补偿、速率 PI、
惯量前馈、电机观测器和差速调度。

\subsection{固件 Rate 模式的摇杆曲线}
\label{sec:rc_curve}

固件 Rate 模式先把 PWM 摇杆归一化到 $r\in[-1,1]$，在
$|r|<0.02$ 时置零，再计算
\begin{equation}
r_e=r\left[e r^2+(1-e)\right],\qquad
\omega_{\mathrm{ref}}=
\operatorname{sat}_{\omega_{\max}}
\left(\frac{200\,R\,r_e}{1-S|r_e|}\right),
\label{eq:rc_curve}
\end{equation}
其中 $R$ 为 \texttt{rc\_rate}，$e$ 为 \texttt{rc\_expo}，$S$ 为
\texttt{rc\_super}。默认值为 $R=0.5$、$e=0.25$，滚转/俯仰的
$S=0.7$，偏航的 $S=0.55$。默认理论满杆速率分别为约
$333\,\mathrm{deg/s}$ 和 $222\,\mathrm{deg/s}$，并受独立的
$600\,\mathrm{deg/s}$ 安全上限约束。该曲线只生成角速度参考，不改变底层执行器分配。

\subsection{限幅、抗饱和与模式切换}

Web 参数源使用
\begin{equation}
|\delta_f|,|\delta_t|\leq25^\circ,qquad
|\Delta\omega|\leq0.7,
\label{eq:actuator_limits}
\end{equation}
固件则将摆角收紧到 $15^\circ$。积分状态在 Web 模式切换时清零，并按参数源钳位；
固件 PI 还使用积分分离和输出限幅。独立钳位不是严格的多输入抗饱和方案：当多个通道同时饱和时，
期望力矩方向可能改变。后续若加入舵面或更多执行器，应采用显式约束控制分配
\cite{johansen2013control}。

\begin{figure}[H]
\centering
\resizebox{0.94\linewidth}{!}{%
\begin{tikzpicture}[>=Stealth,
  node distance=8mm and 5.5mm,
  box/.style={draw=black!65, fill=black!5, rounded corners=2pt, minimum width=22mm, minimum height=14mm, text width=20mm, align=center, font=\footnotesize\sffamily, inner sep=2.5pt},
  smallbox/.style={draw=black!55, fill=black!3, rounded corners=2pt, minimum width=18mm, minimum height=12mm, text width=16mm, align=center, font=\footnotesize\sffamily, inner sep=2pt},
  arrow/.style={->, draw=black!75, line width=.8pt, >={Stealth[length=2.4mm]}},
  feedback/.style={->, draw=blue!65!black, line width=.8pt, >={Stealth[length=2.4mm]}},
  disturbance/.style={->, densely dashed, draw=orange!80!black, line width=.8pt, >={Stealth[length=2.4mm]}},
  group/.style={draw=red!45!black, densely dashed, rounded corners=4pt, fill=red!3, inner sep=4mm},
  annot/.style={font=\scriptsize\sffamily, fill=white, inner sep=1pt},
]
  \node[box, draw=blue!55!black, fill=blue!5] (input) {操纵/上层指令\\$\omega_0,\delta_t$\\$\delta_f,\Delta\omega$};
  \node[box, draw=red!55!black, fill=red!5, minimum width=25mm, text width=23mm, right=of input] (sas) {SAS增稳控制\\角速率阻尼\\姿态$P+I$/速率闭环};
  \node[smallbox, right=of sas] (limit) {物理限幅\\$\delta_{\max}$\\$\Delta\omega_{\max}$};
  \node[box, draw=green!45!black, fill=green!5, minimum width=24mm, text width=22mm, right=of limit] (act) {执行器模型\\电机一阶滞后\\摆角/差速理想瞬时};
  \node[box, minimum width=24mm, text width=22mm, right=of act] (map) {非线性六维\\推力映射\\$\bm F_b,\bm M_b$};
  \node[box, draw=orange!65!black, fill=orange!6, minimum width=24mm, text width=22mm, right=of map] (plant) {六自由度对象\\Newton--Euler\\四元数运动学};

  \draw[arrow] (input) -- (sas);
  \draw[arrow] (sas) -- (limit);
  \draw[arrow] (limit) -- (act);
  \draw[arrow] (act) -- (map);
  \draw[arrow] (map) -- (plant);
  \draw[arrow] (plant.east) -- ++(.8,0) node[annot, right=1pt] {飞行状态};

  \node[smallbox, draw=blue!45!black, fill=blue!4, minimum width=22mm, text width=20mm, below=14mm of plant] (sensor) {理想状态反馈\\$p,q,r,\phi,\theta$};
  \draw[feedback] (plant.south) -- node[annot, right] {真实状态} (sensor.north);
  \draw[feedback] (sensor.west) -- ++(-.65,0) |- node[annot, pos=.32, below] {测量反馈} (sas.south);

  \node[font=\footnotesize\sffamily, text=orange!80!black, above=10mm of plant] (dist) {气动扰动与外力矩};
  \draw[disturbance] (dist) -- (plant.north);

  \begin{scope}[on background layer]
    \node[group, fit=(sas)(sensor), label={[font=\footnotesize\sffamily\bfseries, text=red!55!black]below:SAS闭环}] {};
  \end{scope}
\end{tikzpicture}%
}
\caption{依据Web主循环重绘的级联控制架构。主信号依次经过SAS、物理限幅、执行器模型、六维映射和刚体动力学；当前实现只有电机转速一阶滞后，摆角与差速按理想瞬时指令处理。蓝色回路是当前代码直接使用的理想状态反馈，不代表已实现IMU或状态估计器；其中 $\theta$ 遵循现行显示符号约定。}
\label{fig:control_architecture}
\end{figure}


\subsection{基线的使用边界}

\cref{fig:web_validation}锁定的是 Web 基线在特定初值和参数下的确定性轨迹。
它能够发现符号翻转、参数漂移和积分器行为变化，但不能回答固件闭环是否稳定。
在比较任何新控制器之前，应先统一四元数误差方向、坐标映射、执行器动态、限幅、
工作点和评价指标；当前高保真脚本尚未满足这一条件。
